\par\Needspace{4\baselineskip}
\begin{center}
\hypertarget{native-038-article021}{}\pdfbookmark[2]{21.}{article-021}21.
\end{center}
In the HYPERBOLA, $\varphi$ and $E$ would become imaginary quantities. If we wish to avoid such quantities, we must introduce other auxiliary quantities in their place. We previously denoted the angle whose cosine is $=\frac{1}{e}$ by $\psi$ and found the following expression for the radius vector.
\[
=\frac{p}{2e\cos\frac{1}{2}(v-\psi)\cos\frac{1}{2}(v+\psi)}
\]
The factors in the denominator of this fraction, $\cos\frac{1}{2}(v-\psi)$ and $\cos\frac{1}{2}(v+\psi)$, are equal when $v=0$. The second vanishes at the greatest positive value of $v$, and the first at its most negative value. Thus, if we set $\frac{\cos\frac{1}{2}(v-\psi)}{\cos\frac{1}{2}(v+\psi)}=u$, then $u=1$ at perihelion. It increases without bound as $v$ approaches its limit $180^0-\psi$. Conversely, it keeps decreasing if we suppose that $v$ recedes towards its \SourceLeaf{039}other limit $-(180^0-\psi)$. Thus opposite values of $v$ correspond to reciprocal values of $u$, or, equivalently, to values whose logarithms are negatives of one another.

The quotient $u$ is very convenient as an auxiliary quantity in the hyperbola. The angle whose tangent is $=\operatorname{tang}\frac{1}{2}v\, .\,\BareRoot{\frac{e-1}{e+1}}$ can serve in its place almost as conveniently. Following the analogy with the ellipse, we shall denote this angle by $\frac{1}{2}F$. The following relations among $v$, $r$, $u$, and $F$ are then easily obtained, where we set $a=-b$ so that $b$ is a positive quantity.
\[
\mathrm{I.}\quad b=p\operatorname{cotang}\psi^2
\]
\[
\mathrm{II.}\quad r=\frac{p}{1+e\cos v}=\frac{p\cos\psi}{2\cos\frac{1}{2}(v-\psi)\cos\frac{1}{2}(v+\psi)}
\]
\[
\begin{aligned}
\mathrm{III.}\quad\operatorname{tang}\frac{1}{2}F&=\operatorname{tang}\frac{1}{2}v\, .\,\BareRoot{\frac{e-1}{e+1}}\\
&=\operatorname{tang}\frac{1}{2}v\operatorname{tang}\frac{1}{2}\psi=\frac{u-1}{u+1}
\end{aligned}
\]
\[
\begin{aligned}
\mathrm{IV.}\quad u&=\frac{\cos\frac{1}{2}(v-\psi)}{\cos\frac{1}{2}(v+\psi)}=\frac{1+\operatorname{tang}\frac{1}{2}F}{1-\operatorname{tang}\frac{1}{2}F}\\
&=\operatorname{tang}(45^0+\frac{1}{2}F)
\end{aligned}
\]
\[
\begin{aligned}
\mathrm{V.}\quad\frac{1}{\cos F}&=\frac{1}{2}\left(u+\frac{1}{u}\right)\\
&=\frac{1+\cos\psi\cos v}{2\cos\frac{1}{2}(v-\psi)\cos\frac{1}{2}(v+\psi)}=\frac{e+\cos v}{1+e\cos v}\cdot
\end{aligned}
\]

Subtracting $1$ from both sides of equation V gives
\[
\begin{aligned}
\mathrm{VI.}\quad\sin\frac{1}{2}v\, .\,\BareRoot{r}&=\sin\frac{1}{2}F\, .\,\BareRoot{\frac{p}{(e-1)\cos F}}\\
&=\sin\frac{1}{2}F\, .\,\BareRoot{\frac{(e+1)b}{\cos F}}\\
&=\frac{1}{2}(u-1)\BareRoot{\frac{p}{(e-1)u}}=\frac{1}{2}(u-1)\BareRoot{\frac{(e+1)b}{u}}\cdot
\end{aligned}
\]

Similarly, adding $1$ to both sides gives
\[
\begin{aligned}
\mathrm{VII.}\quad\cos\frac{1}{2}v\, .\,\BareRoot{r}&=\cos\frac{1}{2}F\, .\,\BareRoot{\frac{p}{(e+1)\cos F}}\\
&=\cos\frac{1}{2}F\, .\,\BareRoot{\frac{(e-1)b}{\cos F}}\\
&=\frac{1}{2}(u+1)\BareRoot{\frac{p}{(e+1)u}}=\frac{1}{2}(u+1)\BareRoot{\frac{(e-1)b}{u}}\cdot
\end{aligned}
\]

Dividing VI by VII takes us back to III. Multiplying the two gives
\[
\begin{aligned}
\mathrm{VIII.}\quad r\sin v&=p\operatorname{cotang}\psi\operatorname{tang}F=b\operatorname{tang}\psi\operatorname{tang}F\\
&=\frac{1}{2}p\operatorname{cotang}\psi\, .\,\left(u-\frac{1}{u}\right)=\frac{1}{2}b\operatorname{tang}\psi\, .\,\left(u-\frac{1}{u}\right).
\end{aligned}
\]

Combining equations II and V also readily gives
\[
\mathrm{IX.}\quad r\cos v=b\left(e-\frac{1}{\cos F}\right)=\frac{1}{2}b\left(2e-u-\frac{1}{u}\right)
\]
\[
\mathrm{X.}\quad r=b\left(\frac{e}{\cos F}-1\right)=\frac{1}{2}b\left\{e\left(u+\frac{1}{u}\right)-2\right\}.
\]
